\documentclass[a4paper,11pt]{article}
\usepackage[utf8]{inputenc}
\usepackage[T1]{fontenc}
\usepackage[danish]{babel}
\usepackage{geometry}
\geometry{margin=2.5cm}
\usepackage{hyperref}
\usepackage{enumitem}
\usepackage{xcolor}
\usepackage{array}
\usepackage{pdflscape}
\usepackage{fancyhdr}
\setlength{\headheight}{14pt}
\usepackage{titlesec}

\setlist[itemize]{topsep=2pt,itemsep=2pt,parsep=0pt,leftmargin=1.2cm}

% ----------------------------------------------------------------
% Kildegrundlag (kompileres ikke; dokumentation af afsæt)
%  - Lektionsplan (MD): :contentReference[oaicite:0]{index=0}
%  - Kursusoversigt (MD): :contentReference[oaicite:1]{index=1}
%  - Studieordning, national del (2019): :contentReference[oaicite:2]{index=2}
%  - Studieordning, institutionsdel (2023): :contentReference[oaicite:3]{index=3}
% ----------------------------------------------------------------

% Header og footer
\pagestyle{fancy}
\fancyhf{}
\fancyhead[L]{\textbf{AAMS} - Teknologi \& Projektudvikling}
\fancyhead[R]{\textit{Foreløbig Plan}}
\fancyfoot[C]{\thepage}

% Titel formattering
\titleformat{\section}{\normalfont\Large\bfseries\color{blue!70!black}}{\thesection}{1em}{}

\title{\textbf{\Huge AAMS}\\[0.5cm]
       {\Large Aarhus Maskinmesterskole}\\[1cm]
       {\huge\color{blue!70!black}Detaljeret Forløbsplan}\\[0.3cm]
       {\Large\color{red!70!black}\textit{Teknologi og Projektudvikling — 12 dage (5 ECTS)}}\\[0.5cm]
       {\large\color{gray}\textbf{FORELØBIG VERSION}}}
\author{Underviser: Anders S. Østergaard}
\date{Opdateret: \emph{\today}}

\begin{document}
\maketitle

\tableofcontents
\clearpage

\section*{Formål og kobling til studieordningen}
Forløbet udvikler kompetencer i højniveau-programmering, datakommunikation og teknisk dokumentation samt konfiguration af mindre systemer til kommunikation mellem enheder. Læringsmålene er afledt af fagelementet \emph{Teknologi- og Projektudvikling} (5 ECTS) og indgår i 2. semesterforløbet. Se nationale og institutionelle beskrivelser i studieordningen.

\section*{Læringsmål for forløbet}
Efter forløbet kan den studerende:
\begin{itemize}
  \item Konfigurere datakommunikation mellem controller og eksterne enheder (ESP32, S7/PLC).
  \item Udarbejde mindre Python-programmer til dataopsamling, behandling og visualisering.
  \item Anvende Snap7 til aflæsning af udvalgte DB-variabler fra S7-PLC i Python.
  \item Strukturere dokumentation signal- og blokdiagrammer, testlog mm.
  \item Udarbejde og anvende simple sanity checks, tidsstempling og filtrering af måledata.
  \item Fremlægge løsning, begrunde designvalg og reflektere over datakvalitet og fejlhåndtering.
\end{itemize}

\section*{Evaluering, afleveringer og organisering}
\textbf{Afslutning:} Fremlæggelse på dag 12 med peer-review. \\
\textbf{Software \& hardware:} Python (\texttt{pandas}, \texttt{matplotlib}, \texttt{pyserial}, \texttt{python-snap7}), ESP32 med sensorer, Siemens S7 eller simulator (Snap7).

\clearpage

\section{Dag 1: Teknisk dokumentation og tests}

\textbf{Emne:} Dokumentation, krav og test i projektarbejde.\\
\textbf{Indhold:}
\begin{itemize}
  \item Kravspecifikation (formål, omfang, krav og acceptkriterier).
  \item Teknisk dokumentation (struktur, skabelon, kvalitet).
  \item Tests: testplan, testcases og sporbarhed.
  \item Diagrammer til design/overblik (flowchart, block diagram, state machine).
\end{itemize}
\textbf{Aktivitet:}
\begin{itemize}
  \item Udarbejd en kort kravspecifikation til jeres case.
  \item Start en teknisk dokumentation (evt. ud fra skabelon).
  \item Skriv 3--5 testcases med tydelige acceptkriterier.
\end{itemize}
\textbf{Output:}
\begin{itemize}
  \item Kravspecifikation.
  \item Dokumentationsskitse/skabelon.
  \item Testplan/testcases.
\end{itemize}


\section{Dag 2: Intro til GitHub og Python}

\textbf{Emne:} Versionsstyring og Python-grundforløb.\\
\textbf{Indhold:}
\begin{itemize}
  \item GitHub-workflow: clone, commit, push, branch og merge.
  \item Python intro: køre scripts, datatyper, kontrolstrukturer og fejlfinding.
\end{itemize}
\textbf{Aktivitet:}
\begin{itemize}
  \item Klon/opret et repository og lav små commits.
  \item Arbejd i en branch og merge til \texttt{main}.
  \item Skriv og kør simple Python-programmer.
\end{itemize}
\textbf{Output:}
\begin{itemize}
  \item Repo med historik (commits) og branch/merge.
  \item Python-scripts (grundøvelser).
\end{itemize}


\section{Dag 3: Arbejd videre med dag 02}

\textbf{Emne:} Rutine i GitHub og Python + mini-case.\\
\textbf{Indhold:}
\begin{itemize}
  \item Repetition af GitHub-workflow (små commits, branch/merge).
  \item Repetition/fordybelse i Python (datastrukturer, funktioner, error handling).
  \item Mini-automation case: simulér sensor, log data, timestamp og simple sanity checks.
\end{itemize}
\textbf{Aktivitet:}
\begin{itemize}
  \item Vælg 1--3 fokusområder og arbejd målrettet.
  \item Byg et lille Python-script der logger data struktureret (txt/CSV).
\end{itemize}
\textbf{Output:}
\begin{itemize}
  \item Repo med tydelig historik og en gennemført branch/merge.
  \item Python-script der logger data med timestamp og validering.
\end{itemize}


\section{Dag 4: Python til PLC-kommunikation}

\textbf{Emne:} PLC-kommunikation fra Python (Siemens S7 + Allen-Bradley).\\
\textbf{Indhold:}
\begin{itemize}
  \item Snap7: læs/skriv data til Siemens S7 (BOOL, INT, REAL, WORD).
  \item EtherNet/IP: læs/skriv tags til Allen-Bradley (\texttt{pycomm3}).
  \item Logging (txt/CSV), simple plots/dashboards og fejlhåndtering.
\end{itemize}
\textbf{Aktivitet:}
\begin{itemize}
  \item Implementér læs/skriv for udvalgte variabler/tags.
  \item Log målinger med tidsstempel og visualisér udvalgte signaler.
\end{itemize}
\textbf{Output:}
\begin{itemize}
  \item Python-værktøj til PLC-læs/skriv.
  \item Logfil(er) og simple plots/dashboards.
\end{itemize}


\section{Dag 5: ESP32 intro (MicroPython)}

\textbf{Emne:} ESP32 som sensor-/I/O-enhed og stabilt serielt output.\\
\textbf{Indhold:}
\begin{itemize}
  \item Opsætning af MicroPython på ESP32 (Thonny).
  \item Stabilt output i CSV-lignende format til senere Python-logning.
  \item Afprøvning af sensorer/inputs (fx DHT22, MQ-2, LDR, PIR, distance, PWM, touch).
\end{itemize}
\textbf{Aktivitet:}
\begin{itemize}
  \item Flash MicroPython og verificér forbindelse i Thonny.
  \item Lav serielt output i et stabilt, parsebart format.
  \item Afprøv mindst én sensor og udskriv værdier stabilt.
\end{itemize}
\textbf{Output:}
\begin{itemize}
  \item Kørende ESP32-script og stabil seriel datastrøm.
\end{itemize}


\section{Dag 6: Python + pyserial (seriel datamodtagelse)}

\textbf{Emne:} Seriel kommunikation mellem PC og ESP32.\\
\textbf{Indhold:}
\begin{itemize}
  \item Installation og test af \texttt{pyserial}.
  \item Læsning og parsing af serielle målinger (én linje pr. måling).
  \item Logging til CSV og/eller JSONL med tidsstempel.
  \item (Valgfrit) live-visualisering og simpel realtime processing.
\end{itemize}
\textbf{Aktivitet:}
\begin{itemize}
  \item Byg et Python-script der læser stabilt fra ESP32.
  \item Log data til fil og verificér format og datakvalitet.
\end{itemize}
\textbf{Output:}
\begin{itemize}
  \item Modtager-script + logfil (CSV/JSONL) med timestamp.
  \item (Valgfrit) live-plot og/eller simpel alarmlogik.
\end{itemize}


\section{Dag 7: Python + CSV + Pandas intro}

\textbf{Emne:} Datahåndtering med \texttt{pandas} og \texttt{matplotlib}.\\
\textbf{Indhold:}
\begin{itemize}
  \item Pandas basics: DataFrames.
  \item Import/eksport af CSV.
  \item Datarensning (NaN, outliers, dubletter) og simple tidsserier.
  \item Visualisering af måledata.
\end{itemize}
\textbf{Aktivitet:}
\begin{itemize}
  \item Indlæs et datasæt, rens det, og lav grafer med tydelige akser/titler.
\end{itemize}
\textbf{Output:}
\begin{itemize}
  \item Renset datasæt + Python-script til analyse.
  \item Min. 2 grafer klar til dokumentation.
\end{itemize}


\section{Dag 8: Pandas-visualisering og analyse}

\textbf{Emne:} Professionelle grafer og statistisk bearbejdning.\\
\textbf{Indhold:}
\begin{itemize}
  \item Glidende gennemsnit (\texttt{.rolling()}).
  \item Flere sensorer i én DataFrame og sammenstilling af måleserier.
  \item Annotering og eksport af figurer i høj kvalitet.
  \item Filtrering, outliers og fejlhåndtering.
\end{itemize}
\textbf{Aktivitet:}
\begin{itemize}
  \item Udarbejd mindst én figur med glidende gennemsnit og annotering.
  \item Eksportér figurer til brug i rapport/dokumentation.
\end{itemize}
\textbf{Output:}
\begin{itemize}
  \item Figurer (PNG/PDF) klar til teknisk dokumentation.
\end{itemize}


\section{Dag 9: Sanity checks, tidsstempling og plausibilitet}

\textbf{Emne:} Datakvalitet i sensor- og automationsdata.\\
\textbf{Indhold:}
\begin{itemize}
  \item Sanity checks vs. plausibilitetstests.
  \item Tidsstempling og vurdering af samplingfrekvens.
  \item Outliers, dubletter og fejlformat.
  \item Software-watchdog til overvågning af datakvalitet.
\end{itemize}
\textbf{Aktivitet:}
\begin{itemize}
  \item Implementér en valideringspipeline (sensor \textrightarrow{} check \textrightarrow{} log/visning).
  \item Afprøv watchdog-logik på datastrøm eller testdata.
\end{itemize}
\textbf{Output:}
\begin{itemize}
  \item Sanity/plausibilitet-script og dokumenteret datakvalitet.
\end{itemize}


\section{Dag 10: Dataforståelse og databehandling (Pandas)}

\textbf{Emne:} Fra rå data til renset og valideret datasæt.\\
\textbf{Indhold:}
\begin{itemize}
  \item Datastruktur og datakilder (CSV).
  \item Data inspection (find fejl uden at ændre data).
  \item Data cleaning (NaN, datatyper, dubletter, outliers).
  \item Data validation og mini-rapport (før/efter).
\end{itemize}
\textbf{Aktivitet:}
\begin{itemize}
  \item Generér/brug \texttt{raw\_data.csv} og gennemfør inspection, cleaning og validation.
  \item Skriv en kort rapport med metode og resultater.
\end{itemize}
\textbf{Output:}
\begin{itemize}
  \item Renset/valideret datasæt + mini-rapport.
\end{itemize}


\section{Dag 11: Realtime plotting}

\textbf{Emne:} Live-visualisering af serielle sensordata i Python.\\
\textbf{Indhold:}
\begin{itemize}
  \item Opsætning af miljø (\texttt{matplotlib}, \texttt{pyserial}) og valg af seriel port.
  \item Live plot med \texttt{FuncAnimation}.
  \item Single-sensor plots (temperatur, humidity, LDR, gas, distance).
  \item Multiple plots (samme figur eller subplots) og rullende vindue.
\end{itemize}
\textbf{Aktivitet:}
\begin{itemize}
  \item Byg et realtime plot fra ESP32 og udvid til flere sensorer.
\end{itemize}
\textbf{Output:}
\begin{itemize}
  \item Realtime plotting-script(s) til live data.
\end{itemize}


\section{Dag 12: Opsummering og evaluering}

\textbf{Emne:} Opsamling, refleksion og mini-projekt/fremlæggelse.\\
\textbf{Indhold:}
\begin{itemize}
  \item Opsummering af centrale begreber og workflow (GitHub, ESP32, seriel, data, kvalitet, visualisering).
  \item Evaluering, peer-review og næste skridt mod semesterprojekt.
\end{itemize}
\textbf{Aktivitet:}
\begin{itemize}
  \item Mini-projekt eller fremlæggelse med refleksion over løsning og datakvalitet.
\end{itemize}
\textbf{Output:}
\begin{itemize}
  \item Mini-projekt/fremlæggelse + feedback/opsamling.
\end{itemize}


% ------- OVERSIGTSTABEL -------
\clearpage
\begin{landscape}
\section*{Oversigt pr. dag}
\fontsize{8}{12}\selectfont
\begin{tabular}{|p{0.8cm}|p{1cm}|p{4.2cm}|p{6.2cm}|p{8.2cm}|}
\hline
\textbf{Dag} & \textbf{Lek.} & \textbf{Emne} & \textbf{Indhold} & \textbf{Output \& Materialer} \\
\hline
1 & 1--3 & Teknisk dokumentation \& tests &
Kravspecifikation, dokumentationsstruktur, testcases, acceptkriterier, diagrammer &
Kravspecifikation + testplan/testcases (skitse) \\
\hline
2 & 4--6 & GitHub \& Python intro &
Clone/commit/push, branch/merge, Python-basics og fejlfinding &
Repo med historik + Python-grundøvelser \\
\hline
3 & 7--9 & Arbejd videre med dag 2 &
Repetition GitHub/Python + mini-case (logning, timestamp, sanity checks) &
Repo med branch/merge + log-script \\
\hline
4 & 10--12 & Python til PLC &
Snap7 (S7) og EtherNet/IP (AB): læs/skriv, logging, simple plots &
PLC-værktøj + logfiler/plots \\
\hline
5 & 13--15 & ESP32 intro (MicroPython) &
MicroPython/Thonny, stabilt serielt output, sensorer/inputs &
ESP32-script + stabil datastrøm \\
\hline
6 & 16--18 & Python + pyserial &
Seriel modtagelse/parsing, logging til CSV/JSONL, (valgfrit) live plot &
Modtager-script + logfil(er) \\
\hline
7 & 19--21 & CSV + Pandas intro &
DataFrames, import/eksport, cleaning, time series, visualisering &
Renset datasæt + min. 2 grafer \\
\hline
8 & 22--24 & Pandas visualisering &
Rolling mean, flere sensorer, annotering, eksport, outliers &
Figurer klar til dokumentation \\
\hline
9 & 25--27 & Sanity + timestamp &
Datavalidering, plausibilitetstests, sampling, watchdog &
Validerings-script + datakvalitetsnoter \\
\hline
10 & 28--30 & Databehandling (Pandas) &
Inspection, cleaning, validation og mini-rapport (før/efter) &
Renset/valideret datasæt + rapport \\
\hline
11 & 31--33 & Realtime plotting &
Live plot med \texttt{FuncAnimation}, single/multiple sensor plots &
Realtime plotting-script(s) \\
\hline
12 & 34--36 & Opsummering \& evaluering &
Opsamling, refleksion, mini-projekt/fremlæggelse, peer-review &
Mini-projekt/fremlæggelse + feedback \\
\hline
\end{tabular}
\end{landscape}


\end{document}
